\chapter{La géométrie : profils NACA à quatre chiffres}\label{ch:naca}
\iterbanner{V0}

\section{Trois paramètres, deux lois}
Un profil NACA $MPXX$ superpose une \textbf{ligne moyenne} (la cambrure) et une \textbf{loi d'épaisseur} symétrique. Trois nombres le définissent, et ce sont les trois variables de conception du hackathon :
\begin{itemize}
\item $m=M/100$ : la cambrure maximale, en fraction de corde ;
\item $p=P/10$ : l'abscisse de cette cambrure maximale ;
\item $t=XX/100$ : l'épaisseur relative maximale.
\end{itemize}
Le NACA 2412 a donc $m=0{,}02$, $p=0{,}4$ et $t=0{,}12$.

\paragraph{La ligne moyenne.} Elle est formée de deux arcs de parabole qui se raccordent en $x=p$, au sommet de la cambrure :
\[
y_c(x)=\begin{cases}\dfrac{m}{p^2}\big(2px-x^2\big) & 0\le x<p,\\[8pt]
\dfrac{m}{(1-p)^2}\big((1-2p)+2px-x^2\big) & p\le x\le1,\end{cases}
\qquad
\frac{\dd y_c}{\dd x}=\begin{cases}\dfrac{2m}{p^2}(p-x) & x<p,\\[8pt] \dfrac{2m}{(1-p)^2}(p-x) & x\ge p.\end{cases}
\]
Les deux arcs valent $m$ en $x=p$, avec une pente nulle : la ligne est continue et dérivable. Elle s'annule aux deux bords ($y_c(0)=y_c(1)=0$). Sa courbure, en revanche, saute en $x=p$, de $-2m/p^2$ à $-2m/(1-p)^2$.

\paragraph{La loi d'épaisseur.} C'est un polynôme calé sur des profils de bonne qualité dans les années 1930 :
\[
y_t(x)=5t\big(0{,}2969\sqrt x-0{,}1260\,x-0{,}3516\,x^2+0{,}2843\,x^3+a_4\,x^4\big).
\]
Le coefficient d'origine $a_4=-0{,}1015$ laisse un bord de fuite d'épaisseur finie ($y_t(1)=0{,}0021\cdot5t$). Le code utilise $a_4=-0{,}1036$, qui donne $y_t(1)=0$ : un \textbf{bord de fuite fermé}. C'est nécessaire pour un contour de panneaux fermé et pour un STL étanche. Le terme en $\sqrt x$ donne un bord d'attaque arrondi, de rayon
\[
r_{BA}=1{,}1019\,t^2 .
\]
L'épaisseur maximale vaut $t$ et se trouve vers $x\approx0{,}30$, \emph{quelle que soit} la cambrure.

\paragraph{L'assemblage.} L'épaisseur est appliquée \emph{perpendiculairement} à la ligne moyenne (figure~\ref{fig:naca-construction}). Avec $\vartheta=\arctan(\dd y_c/\dd x)$ :
\begin{align*}
x_{u}&=x-y_t\sin\vartheta, & y_{u}&=y_c+y_t\cos\vartheta & &\text{(extrados)},\\
x_{l}&=x+y_t\sin\vartheta, & y_{l}&=y_c-y_t\cos\vartheta & &\text{(intrados)}.
\end{align*}

\begin{figure}[htbp]
\centering
\begin{tikzpicture}[x=10.5cm,y=10.5cm]
\draw[cmuted,thin] (0,0) -- (1,0) node[right,font=\small] {corde};
\draw[cblue,thick,dashed] plot[domain=0:0.4,samples=30] (\x,{0.05/0.16*(0.8*\x-\x*\x)}) -- plot[domain=0.4:1,samples=40] (\x,{0.05/0.36*(0.2+0.8*\x-\x*\x)});
\node[cblue,font=\small] at (0.75,0.055) {ligne moyenne $y_c$};
% point de construction
\def\xx{0.22}
\pgfmathsetmacro{\yc}{0.05/0.16*(0.8*\xx-\xx*\xx)}
\pgfmathsetmacro{\sl}{2*0.05/0.16*(0.4-\xx)}
\pgfmathsetmacro{\th}{atan(\sl)}
\pgfmathsetmacro{\yt}{5*0.12*(0.2969*sqrt(\xx)-0.1260*\xx-0.3516*\xx*\xx+0.2843*\xx*\xx*\xx-0.1036*\xx*\xx*\xx*\xx)}
\fill[cink] (\xx,\yc) circle (0.2pt);
\draw[corange,very thick,-{Stealth[length=2mm]}] (\xx,\yc) -- ++(90+\th:\yt) node[above left,font=\small] {$y_t$};
\draw[corange,very thick,-{Stealth[length=2mm]}] (\xx,\yc) -- ++(-90+\th:\yt) node[below right,font=\small] {$y_t$};
\draw[cink,thin] (\xx,\yc) -- ++(\th:0.12);
\draw[cink,thin] (\xx,\yc) -- ++(0:0.12);
\draw[cink] (\xx+0.07,\yc) arc (0:\th:0.07);
\node[font=\small] at (\xx+0.105,\yc+0.012) {$\vartheta$};
% profil complet
\draw[cink,thick] plot[smooth,tension=0.5] coordinates {(1,0) (0.8,0.032) (0.6,0.058) (0.4,0.077) (0.25,0.08) (0.12,0.07) (0.04,0.042) (0,0) (0.04,-0.026) (0.12,-0.036) (0.25,-0.039) (0.4,-0.034) (0.6,-0.022) (0.8,-0.011) (1,0)};
\draw[cmuted,{Stealth[length=1.5mm]}-{Stealth[length=1.5mm]}] (0,-0.075) -- (0.4,-0.075) node[midway,below,font=\small] {$p$};
\draw[cmuted,{Stealth[length=1.5mm]}-{Stealth[length=1.5mm]}] (0.4,0) -- (0.4,0.05) node[midway,right,font=\small] {$m$};
\end{tikzpicture}
\caption{Construction d'un profil NACA (cambrure exagérée) : l'épaisseur $y_t$ est reportée de part et d'autre de la ligne moyenne, perpendiculairement à celle-ci.}\label{fig:naca-construction}
\end{figure}

\section{Ce que chaque paramètre fait physiquement}
\begin{figure}[htbp]
\centering
% figure générée par tools/make_figures.py : parametres
\begin{tikzpicture}
\begin{groupplot}[group style={group size=1 by 3,vertical sep=1.5cm,x descriptions at=edge bottom},width=\linewidth,height=3.9cm,
  axis equal image,xmin=-0.02,xmax=1.02,ymin=-0.12,ymax=0.17,xlabel={$x/c$},ylabel={$y/c$},
  title style={font=\small},legend columns=3,legend style={draw=none,font=\scriptsize,at={(0.99,0.97)},anchor=north east,fill=white,fill opacity=0.8,text opacity=1}]
\nextgroupplot[title={effet de $m$ ($p=0{,}4$, $t=12$~\%)}]
\addplot[cmuted,thick] coordinates {(1,-1.6653e-17) (0.99931,9.9554e-05) (0.99726,0.00039741) (0.99384,0.00089119) (0.98907,0.0015769) (0.98296,0.0024493) (0.97553,0.0035014) (0.96679,0.0047252) (0.95677,0.0061118) (0.9455,0.0076508) (0.93301,0.0093315) (0.91934,0.011142) (0.90451,0.013071) (0.88857,0.015105) (0.87157,0.017232) (0.85355,0.019438) (0.83457,0.021712) (0.81466,0.024038) (0.79389,0.026405) (0.77232,0.028798) (0.75,0.031204) (0.727,0.03361) (0.70337,0.036002) (0.67918,0.038366) (0.65451,0.040686) (0.62941,0.042949) (0.60396,0.04514) (0.57822,0.047242) (0.55226,0.049241) (0.52617,0.051119) (0.5,0.052862) (0.47383,0.054451) (0.44774,0.055872) (0.42178,0.057108) (0.39604,0.058144) (0.37059,0.058965) (0.34549,0.059557) (0.32082,0.059907) (0.29663,0.060005) (0.273,0.059841) (0.25,0.059407) (0.22768,0.058699) (0.20611,0.057712) (0.18534,0.056446) (0.16543,0.054901) (0.14645,0.053083) (0.12843,0.050995) (0.11143,0.048648) (0.095492,0.046049) (0.080665,0.04321) (0.066987,0.040145) (0.054497,0.036867) (0.043227,0.033389) (0.03321,0.029726) (0.024472,0.025893) (0.017037,0.021904) (0.010926,0.01777) (0.0061558,0.013503) (0.0027391,0.0091145) (0.00068523,0.0046113) (0,0) (0.00068523,-0.0046113) (0.0027391,-0.0091145) (0.0061558,-0.013503) (0.010926,-0.01777) (0.017037,-0.021904) (0.024472,-0.025893) (0.03321,-0.029726) (0.043227,-0.033389) (0.054497,-0.036867) (0.066987,-0.040145) (0.080665,-0.04321) (0.095492,-0.046049) (0.11143,-0.048648) (0.12843,-0.050995) (0.14645,-0.053083) (0.16543,-0.054901) (0.18534,-0.056446) (0.20611,-0.057712) (0.22768,-0.058699) (0.25,-0.059407) (0.273,-0.059841) (0.29663,-0.060005) (0.32082,-0.059907) (0.34549,-0.059557) (0.37059,-0.058965) (0.39604,-0.058144) (0.42178,-0.057108) (0.44774,-0.055872) (0.47383,-0.054451) (0.5,-0.052862) (0.52617,-0.051119) (0.55226,-0.049241) (0.57822,-0.047242) (0.60396,-0.04514) (0.62941,-0.042949) (0.65451,-0.040686) (0.67918,-0.038366) (0.70337,-0.036002) (0.727,-0.03361) (0.75,-0.031204) (0.77232,-0.028798) (0.79389,-0.026405) (0.81466,-0.024038) (0.83457,-0.021712) (0.85355,-0.019438) (0.87157,-0.017232) (0.88857,-0.015105) (0.90451,-0.013071) (0.91934,-0.011142) (0.93301,-0.0093315) (0.9455,-0.0076508) (0.95677,-0.0061118) (0.96679,-0.0047252) (0.97553,-0.0035014) (0.98296,-0.0024493) (0.98907,-0.0015769) (0.99384,-0.00089119) (0.99726,-0.00039741) (0.99931,-9.9554e-05) (1,1.6653e-17)};\addlegendentry{$m=0$}
\addplot[cmuted,dashed,thin,forget plot] coordinates {(0,0) (0.016949,0) (0.033898,0) (0.050847,0) (0.067797,0) (0.084746,0) (0.10169,0) (0.11864,0) (0.13559,0) (0.15254,0) (0.16949,0) (0.18644,0) (0.20339,0) (0.22034,0) (0.23729,0) (0.25424,0) (0.27119,0) (0.28814,0) (0.30508,0) (0.32203,0) (0.33898,0) (0.35593,0) (0.37288,0) (0.38983,0) (0.40678,0) (0.42373,0) (0.44068,0) (0.45763,0) (0.47458,0) (0.49153,0) (0.50847,0) (0.52542,0) (0.54237,0) (0.55932,0) (0.57627,0) (0.59322,0) (0.61017,0) (0.62712,0) (0.64407,0) (0.66102,0) (0.67797,0) (0.69492,0) (0.71186,0) (0.72881,0) (0.74576,0) (0.76271,0) (0.77966,0) (0.79661,0) (0.81356,0) (0.83051,0) (0.84746,0) (0.86441,0) (0.88136,0) (0.89831,0) (0.91525,0) (0.9322,0) (0.94915,0) (0.9661,0) (0.98305,0) (1,0)};
\addplot[cblue,very thick] coordinates {(1,-1.6571e-17) (0.99932,0.00016755) (0.9973,0.00066874) (0.99393,0.0014993) (0.98923,0.0026521) (0.9832,0.0041173) (0.97586,0.0058826) (0.96723,0.0079334) (0.95734,0.010253) (0.9462,0.012822) (0.93384,0.01562) (0.9203,0.018625) (0.9056,0.021814) (0.8898,0.025163) (0.87292,0.028647) (0.85502,0.032241) (0.83613,0.035918) (0.81632,0.039652) (0.79562,0.043419) (0.7741,0.047191) (0.75182,0.050943) (0.72882,0.05465) (0.70519,0.058287) (0.68097,0.061829) (0.65623,0.065252) (0.63105,0.068532) (0.60549,0.071647) (0.57962,0.074575) (0.55351,0.077293) (0.52724,0.079781) (0.50088,0.082021) (0.4745,0.083993) (0.44818,0.08568) (0.42199,0.087068) (0.39596,0.088141) (0.36994,0.088799) (0.34427,0.088987) (0.31904,0.088705) (0.29431,0.087957) (0.27016,0.086749) (0.24666,0.085095) (0.2239,0.083009) (0.20192,0.080511) (0.18081,0.077624) (0.16062,0.074374) (0.14142,0.07079) (0.12326,0.066905) (0.10619,0.062751) (0.090267,0.058365) (0.075527,0.053784) (0.062013,0.049042) (0.04976,0.044179) (0.0388,0.039228) (0.029159,0.034224) (0.020861,0.029199) (0.013923,0.024182) (0.0083607,0.0192) (0.0041829,0.014275) (0.0013961,0.0094244) (2.3425e-06,0.0046631) (0,0) (0.0013681,-0.0044577) (0.004082,-0.0086055) (0.0081288,-0.012442) (0.013492,-0.015967) (0.020151,-0.01918) (0.028082,-0.022082) (0.03726,-0.024674) (0.047655,-0.02696) (0.059234,-0.028943) (0.071962,-0.030629) (0.085802,-0.032024) (0.10072,-0.033138) (0.11666,-0.033979) (0.13359,-0.034561) (0.15147,-0.034899) (0.17025,-0.035007) (0.18987,-0.034903) (0.21029,-0.034609) (0.23147,-0.034144) (0.25334,-0.033532) (0.27585,-0.032797) (0.29896,-0.031964) (0.32259,-0.031057) (0.34671,-0.030101) (0.37124,-0.029123) (0.39613,-0.028147) (0.42158,-0.027147) (0.44729,-0.02606) (0.47316,-0.024901) (0.49912,-0.023687) (0.52509,-0.022434) (0.55102,-0.021157) (0.57681,-0.019868) (0.60242,-0.01858) (0.62777,-0.017304) (0.65278,-0.016048) (0.6774,-0.014819) (0.70155,-0.013626) (0.72517,-0.012471) (0.74818,-0.01136) (0.77054,-0.010294) (0.79216,-0.0092772) (0.813,-0.0083095) (0.833,-0.0073921) (0.85209,-0.0065257) (0.87022,-0.0057107) (0.88735,-0.0049472) (0.90341,-0.0042357) (0.91837,-0.0035765) (0.93219,-0.0029701) (0.94481,-0.0024172) (0.95621,-0.0019186) (0.96635,-0.0014752) (0.97519,-0.0010881) (0.98273,-0.00075825) (0.98892,-0.00048672) (0.99376,-0.00027443) (0.99722,-0.00012218) (0.9993,-3.0577e-05) (1,1.6571e-17)};\addlegendentry{$m=3~\%$}
\addplot[cblue,dashed,thin,forget plot] coordinates {(0,0) (0.016949,0.0024885) (0.033898,0.0048693) (0.050847,0.0071423) (0.067797,0.0093077) (0.084746,0.011365) (0.10169,0.013315) (0.11864,0.015157) (0.13559,0.016892) (0.15254,0.018518) (0.16949,0.020037) (0.18644,0.021449) (0.20339,0.022752) (0.22034,0.023948) (0.23729,0.025036) (0.25424,0.026016) (0.27119,0.026889) (0.28814,0.027654) (0.30508,0.028311) (0.32203,0.02886) (0.33898,0.029302) (0.35593,0.029636) (0.37288,0.029862) (0.38983,0.029981) (0.40678,0.029996) (0.42373,0.029953) (0.44068,0.029862) (0.45763,0.029723) (0.47458,0.029537) (0.49153,0.029302) (0.50847,0.029019) (0.52542,0.028689) (0.54237,0.028311) (0.55932,0.027885) (0.57627,0.027411) (0.59322,0.026889) (0.61017,0.026319) (0.62712,0.025701) (0.64407,0.025036) (0.66102,0.024323) (0.67797,0.023561) (0.69492,0.022752) (0.71186,0.021895) (0.72881,0.02099) (0.74576,0.020037) (0.76271,0.019037) (0.77966,0.017988) (0.79661,0.016892) (0.81356,0.015747) (0.83051,0.014555) (0.84746,0.013315) (0.86441,0.012027) (0.88136,0.010691) (0.89831,0.0093077) (0.91525,0.0078761) (0.9322,0.0063966) (0.94915,0.0048693) (0.9661,0.0032941) (0.98305,0.001671) (1,0)};
\addplot[corange,very thick] coordinates {(1,-1.633e-17) (0.99933,0.00023459) (0.99734,0.00093632) (0.99402,0.0020991) (0.98938,0.0037127) (0.98343,0.0057633) (0.97619,0.0082332) (0.96767,0.011101) (0.95789,0.014343) (0.94687,0.017932) (0.93465,0.021837) (0.92124,0.026027) (0.90668,0.030468) (0.891,0.035125) (0.87425,0.039959) (0.85646,0.044935) (0.83768,0.050013) (0.81795,0.055155) (0.79733,0.060321) (0.77587,0.065475) (0.75362,0.070578) (0.73064,0.075592) (0.70699,0.080481) (0.68274,0.08521) (0.65795,0.089745) (0.63268,0.094053) (0.60702,0.098103) (0.58102,0.10187) (0.55476,0.10531) (0.52832,0.10842) (0.50176,0.11117) (0.47517,0.11353) (0.44862,0.11549) (0.4222,0.11703) (0.39587,0.11814) (0.36929,0.11863) (0.34306,0.11839) (0.31726,0.11745) (0.29199,0.11582) (0.26733,0.11352) (0.24336,0.1106) (0.22016,0.10708) (0.1978,0.10301) (0.17637,0.098448) (0.15592,0.093438) (0.13653,0.08804) (0.11825,0.082313) (0.10114,0.076319) (0.085239,0.070121) (0.0706,0.063781) (0.057259,0.057362) (0.045249,0.050923) (0.034597,0.044522) (0.025325,0.038211) (0.017452,0.032041) (0.01099,0.026055) (0.0059485,0.020291) (0.0023306,0.014783) (0.00013649,0.0095539) (-0.00063772,0.0046228) (0,0) (0.0020082,-0.004212) (0.0053416,-0.0079161) (0.0099811,-0.011118) (0.015904,-0.013825) (0.023084,-0.01605) (0.031491,-0.017807) (0.041094,-0.019112) (0.051858,-0.019987) (0.063744,-0.020453) (0.076715,-0.020535) (0.090729,-0.020263) (0.10574,-0.019665) (0.12172,-0.018775) (0.13861,-0.017626) (0.15636,-0.016257) (0.17495,-0.014704) (0.19431,-0.013008) (0.21441,-0.011209) (0.2352,-0.0093499) (0.25664,-0.0074726) (0.27868,-0.0056194) (0.30127,-0.0038326) (0.32437,-0.0021532) (0.34792,-0.00062125) (0.37189,0.00072516) (0.39622,0.0018503) (0.42137,0.0028142) (0.44685,0.0037552) (0.47249,0.0046567) (0.49824,0.0055012) (0.52402,0.0062729) (0.54977,0.0069586) (0.57542,0.0075474) (0.60089,0.0080312) (0.62613,0.0084043) (0.65107,0.0086637) (0.67563,0.0088089) (0.69975,0.0088418) (0.72335,0.0087665) (0.74638,0.0085892) (0.76877,0.0083177) (0.79046,0.0079615) (0.81137,0.0075311) (0.83145,0.0070382) (0.85065,0.0064948) (0.8689,0.0059137) (0.88614,0.0053074) (0.90234,0.0046886) (0.91743,0.0040695) (0.93138,0.003462) (0.94413,0.002877) (0.95566,0.0023249) (0.96591,0.001815) (0.97487,0.0013559) (0.9825,0.00095476) (0.98877,0.00061796) (0.99367,0.00035063) (0.99718,0.0001568) (0.9993,3.9343e-05) (1,1.633e-17)};\addlegendentry{$m=6~\%$}
\addplot[corange,dashed,thin,forget plot] coordinates {(0,0) (0.016949,0.004977) (0.033898,0.0097386) (0.050847,0.014285) (0.067797,0.018615) (0.084746,0.022731) (0.10169,0.02663) (0.11864,0.030315) (0.13559,0.033783) (0.15254,0.037037) (0.16949,0.040075) (0.18644,0.042897) (0.20339,0.045504) (0.22034,0.047896) (0.23729,0.050072) (0.25424,0.052032) (0.27119,0.053778) (0.28814,0.055307) (0.30508,0.056622) (0.32203,0.05772) (0.33898,0.058604) (0.35593,0.059272) (0.37288,0.059724) (0.38983,0.059961) (0.40678,0.059992) (0.42373,0.059906) (0.44068,0.059724) (0.45763,0.059447) (0.47458,0.059073) (0.49153,0.058604) (0.50847,0.058039) (0.52542,0.057378) (0.54237,0.056622) (0.55932,0.055769) (0.57627,0.054821) (0.59322,0.053778) (0.61017,0.052638) (0.62712,0.051403) (0.64407,0.050072) (0.66102,0.048645) (0.67797,0.047122) (0.69492,0.045504) (0.71186,0.04379) (0.72881,0.04198) (0.74576,0.040075) (0.76271,0.038073) (0.77966,0.035976) (0.79661,0.033783) (0.81356,0.031495) (0.83051,0.02911) (0.84746,0.02663) (0.86441,0.024054) (0.88136,0.021383) (0.89831,0.018615) (0.91525,0.015752) (0.9322,0.012793) (0.94915,0.0097386) (0.9661,0.0065881) (0.98305,0.003342) (1,0)};
\nextgroupplot[title={effet de $p$ ($m=5$~\%, $t=12$~\%)}]
\addplot[cmuted,thick] coordinates {(1,-1.6525e-17) (0.99933,0.0001844) (0.99731,0.00073616) (0.99395,0.0016509) (0.98927,0.0029215) (0.98326,0.0045381) (0.97595,0.0064881) (0.96735,0.0087568) (0.95749,0.011327) (0.94639,0.01418) (0.93407,0.017294) (0.92058,0.020647) (0.90594,0.024217) (0.89019,0.027977) (0.87337,0.031903) (0.85553,0.035968) (0.83671,0.040147) (0.81696,0.044412) (0.79633,0.048736) (0.77488,0.053094) (0.75267,0.057457) (0.72975,0.0618) (0.70619,0.066096) (0.68205,0.07032) (0.65739,0.074445) (0.63228,0.078447) (0.6068,0.082302) (0.581,0.085984) (0.55497,0.089472) (0.52877,0.092742) (0.50248,0.095772) (0.47616,0.098543) (0.4499,0.10104) (0.42376,0.10323) (0.39782,0.10511) (0.37216,0.10667) (0.34685,0.10789) (0.32195,0.10876) (0.29754,0.10927) (0.27369,0.10942) (0.25046,0.10921) (0.22793,0.10864) (0.20616,0.10771) (0.18327,0.10614) (0.16071,0.1032) (0.1394,0.099028) (0.11945,0.093795) (0.10091,0.087691) (0.083851,0.080901) (0.068311,0.073606) (0.05432,0.065979) (0.041894,0.058182) (0.031044,0.050364) (0.021769,0.042663) (0.014067,0.035198) (0.0079261,0.028074) (0.0033323,0.021379) (0.00026701,0.015182) (-0.0012922,0.0095347) (-0.0013713,0.0044693) (0,0) (0.0027418,-0.0037852) (0.0067703,-0.0068144) (0.012045,-0.0091211) (0.01852,-0.010752) (0.026148,-0.011763) (0.034877,-0.012224) (0.04465,-0.01221) (0.055411,-0.011809) (0.067099,-0.01111) (0.079655,-0.01021) (0.093018,-0.0092081) (0.10713,-0.0082058) (0.12194,-0.0073036) (0.13741,-0.0066015) (0.15349,-0.0061982) (0.17016,-0.0061909) (0.18741,-0.0066763) (0.20605,-0.0077147) (0.22743,-0.0087581) (0.24954,-0.009601) (0.27232,-0.010254) (0.29573,-0.010728) (0.31969,-0.011037) (0.34414,-0.011195) (0.36902,-0.011217) (0.39426,-0.011119) (0.4198,-0.010917) (0.44557,-0.010625) (0.4715,-0.01026) (0.49752,-0.0098348) (0.52357,-0.0093643) (0.54956,-0.0088608) (0.57543,-0.0083355) (0.60111,-0.0077986) (0.62653,-0.0072586) (0.65163,-0.0067229) (0.67632,-0.0061973) (0.70055,-0.0056866) (0.72424,-0.0051942) (0.74733,-0.0047226) (0.76975,-0.0042733) (0.79145,-0.003847) (0.81236,-0.0034442) (0.83242,-0.0030645) (0.85158,-0.0027076) (0.86977,-0.002373) (0.88696,-0.0020601) (0.90308,-0.0017684) (0.91809,-0.0014978) (0.93195,-0.0012481) (0.94462,-0.0010194) (0.95606,-0.00081205) (0.96623,-0.00062663) (0.97511,-0.00046375) (0.98267,-0.00032417) (0.98888,-0.00020864) (0.99373,-0.00011789) (0.99721,-5.257e-05) (0.9993,-1.3169e-05) (1,1.6525e-17)};\addlegendentry{$p=0{,}20$}
\addplot[cmuted,dashed,thin,forget plot] coordinates {(0,0) (0.016949,0.0081155) (0.033898,0.015513) (0.050847,0.022192) (0.067797,0.028153) (0.084746,0.033396) (0.10169,0.03792) (0.11864,0.041727) (0.13559,0.044815) (0.15254,0.047185) (0.16949,0.048837) (0.18644,0.04977) (0.20339,0.049999) (0.22034,0.049968) (0.23729,0.049891) (0.25424,0.04977) (0.27119,0.049604) (0.28814,0.049393) (0.30508,0.049137) (0.32203,0.048837) (0.33898,0.048491) (0.35593,0.0481) (0.37288,0.047665) (0.38983,0.047185) (0.40678,0.04666) (0.42373,0.046089) (0.44068,0.045475) (0.45763,0.044815) (0.47458,0.04411) (0.49153,0.04336) (0.50847,0.042566) (0.52542,0.041727) (0.54237,0.040842) (0.55932,0.039913) (0.57627,0.038939) (0.59322,0.03792) (0.61017,0.036856) (0.62712,0.035748) (0.64407,0.034594) (0.66102,0.033396) (0.67797,0.032152) (0.69492,0.030864) (0.71186,0.029531) (0.72881,0.028153) (0.74576,0.02673) (0.76271,0.025262) (0.77966,0.023749) (0.79661,0.022192) (0.81356,0.020589) (0.83051,0.018942) (0.84746,0.01725) (0.86441,0.015513) (0.88136,0.013731) (0.89831,0.011904) (0.91525,0.010032) (0.9322,0.0081155) (0.94915,0.0061539) (0.9661,0.0041475) (0.98305,0.0020962) (1,0)};
\addplot[cblue,very thick] coordinates {(1,-1.6427e-17) (0.99933,0.00021234) (0.99733,0.00084752) (0.99399,0.0019) (0.98933,0.0033607) (0.98335,0.005217) (0.97608,0.0074529) (0.96753,0.01005) (0.95771,0.012985) (0.94665,0.016235) (0.93438,0.019772) (0.92093,0.023568) (0.90632,0.027593) (0.8906,0.031815) (0.87381,0.0362) (0.85598,0.040715) (0.83717,0.045326) (0.81741,0.049999) (0.79676,0.054699) (0.77528,0.059392) (0.75302,0.064044) (0.73004,0.068622) (0.70639,0.073093) (0.68215,0.077425) (0.65738,0.081588) (0.63214,0.085553) (0.60651,0.08929) (0.58055,0.092773) (0.55435,0.095977) (0.52796,0.098877) (0.50147,0.10145) (0.47495,0.10368) (0.44848,0.10555) (0.42213,0.10704) (0.3959,0.10814) (0.36951,0.10868) (0.34346,0.10859) (0.31785,0.10787) (0.29276,0.10654) (0.26827,0.10461) (0.24445,0.10212) (0.22139,0.099082) (0.19916,0.095544) (0.17783,0.091545) (0.15747,0.087127) (0.13814,0.082338) (0.11989,0.077229) (0.10279,0.071852) (0.086882,0.06626) (0.072207,0.060508) (0.058807,0.054647) (0.046715,0.048732) (0.035961,0.042811) (0.026568,0.036932) (0.018555,0.031139) (0.011938,0.02547) (0.0067275,0.019961) (0.0029283,0.014639) (0.00054272,0.0095283) (-0.00043136,0.0046452) (0,0) (0.0018018,-0.0043029) (0.0049354,-0.0081635) (0.0093834,-0.011585) (0.015125,-0.014572) (0.022136,-0.017133) (0.030388,-0.019278) (0.039852,-0.021017) (0.050494,-0.022366) (0.062278,-0.02334) (0.075168,-0.023958) (0.089122,-0.024242) (0.1041,-0.024214) (0.12006,-0.023898) (0.13696,-0.023324) (0.15475,-0.022519) (0.1734,-0.021515) (0.19285,-0.020344) (0.21305,-0.019041) (0.23397,-0.017641) (0.25555,-0.016179) (0.27774,-0.014693) (0.3005,-0.013219) (0.32378,-0.011793) (0.34752,-0.010451) (0.37167,-0.0092252) (0.39619,-0.0081488) (0.42144,-0.0071731) (0.44699,-0.0061837) (0.47272,-0.0051969) (0.49853,-0.00423) (0.52438,-0.0032987) (0.55018,-0.0024167) (0.57588,-0.0015957) (0.6014,-0.00084498) (0.62668,-0.00017188) (0.65164,0.00041864) (0.67622,0.00092376) (0.70035,0.0013427) (0.72395,0.0016766) (0.74698,0.0019282) (0.76936,0.0021018) (0.79102,0.0022032) (0.81191,0.0022388) (0.83196,0.0022161) (0.85112,0.0021431) (0.86933,0.0020278) (0.88654,0.0018788) (0.90269,0.0017043) (0.91774,0.0015123) (0.93165,0.0013104) (0.94436,0.0011058) (0.95584,0.00090506) (0.96606,0.00071404) (0.97498,0.00053799) (0.98257,0.00038144) (0.98882,0.00024822) (0.9937,0.00014141) (0.9972,6.3412e-05) (0.9993,1.5937e-05) (1,1.6427e-17)};\addlegendentry{$p=0{,}40$}
\addplot[cblue,dashed,thin,forget plot] coordinates {(0,0) (0.016949,0.0041475) (0.033898,0.0081155) (0.050847,0.011904) (0.067797,0.015513) (0.084746,0.018942) (0.10169,0.022192) (0.11864,0.025262) (0.13559,0.028153) (0.15254,0.030864) (0.16949,0.033396) (0.18644,0.035748) (0.20339,0.03792) (0.22034,0.039913) (0.23729,0.041727) (0.25424,0.04336) (0.27119,0.044815) (0.28814,0.046089) (0.30508,0.047185) (0.32203,0.0481) (0.33898,0.048837) (0.35593,0.049393) (0.37288,0.04977) (0.38983,0.049968) (0.40678,0.049994) (0.42373,0.049922) (0.44068,0.04977) (0.45763,0.049539) (0.47458,0.049228) (0.49153,0.048837) (0.50847,0.048366) (0.52542,0.047815) (0.54237,0.047185) (0.55932,0.046475) (0.57627,0.045685) (0.59322,0.044815) (0.61017,0.043865) (0.62712,0.042836) (0.64407,0.041727) (0.66102,0.040538) (0.67797,0.039269) (0.69492,0.03792) (0.71186,0.036492) (0.72881,0.034984) (0.74576,0.033396) (0.76271,0.031728) (0.77966,0.02998) (0.79661,0.028153) (0.81356,0.026246) (0.83051,0.024259) (0.84746,0.022192) (0.86441,0.020045) (0.88136,0.017819) (0.89831,0.015513) (0.91525,0.013127) (0.9322,0.010661) (0.94915,0.0081155) (0.9661,0.0054901) (0.98305,0.002785) (1,0)};
\addplot[corange,very thick] coordinates {(1,-1.6257e-17) (0.99934,0.00024935) (0.99735,0.00099489) (0.99403,0.0022291) (0.98941,0.0039397) (0.98348,0.0061095) (0.97625,0.0087168) (0.96774,0.011736) (0.95798,0.015137) (0.94697,0.018886) (0.93475,0.022947) (0.92133,0.02728) (0.90676,0.031844) (0.89106,0.036594) (0.87428,0.041486) (0.85644,0.046472) (0.83759,0.051506) (0.81778,0.05654) (0.79705,0.061528) (0.77546,0.066422) (0.75307,0.071177) (0.72992,0.075748) (0.70609,0.080092) (0.68163,0.084167) (0.65661,0.087935) (0.63109,0.091359) (0.60516,0.094405) (0.57888,0.097041) (0.55232,0.099239) (0.52577,0.10102) (0.49913,0.10244) (0.47246,0.10348) (0.44585,0.10411) (0.41936,0.10434) (0.39309,0.10415) (0.3671,0.10354) (0.34147,0.10251) (0.31629,0.10105) (0.29162,0.099185) (0.26755,0.09691) (0.24414,0.094241) (0.22146,0.091197) (0.19959,0.087795) (0.17858,0.08406) (0.15851,0.080018) (0.13943,0.075698) (0.12139,0.071132) (0.10445,0.066351) (0.088649,0.061393) (0.07404,0.05629) (0.060657,0.051081) (0.048537,0.045799) (0.037711,0.040481) (0.028204,0.035158) (0.02004,0.029862) (0.013237,0.024621) (0.0078086,0.019461) (0.0037665,0.014403) (0.0011165,0.0094657) (-0.00013866,0.0046616) (0,0) (0.0015091,-0.0044125) (0.0043616,-0.0084721) (0.0085452,-0.012177) (0.014044,-0.015527) (0.020838,-0.018522) (0.028904,-0.021161) (0.038216,-0.023446) (0.048744,-0.025379) (0.060456,-0.026964) (0.073317,-0.028205) (0.08729,-0.029109) (0.10233,-0.029683) (0.11841,-0.029937) (0.13547,-0.029883) (0.15347,-0.029535) (0.17236,-0.028908) (0.1921,-0.02802) (0.21263,-0.02689) (0.2339,-0.02554) (0.25586,-0.023994) (0.27846,-0.022274) (0.30164,-0.020407) (0.32534,-0.018418) (0.34951,-0.016334) (0.37408,-0.014182) (0.399,-0.011987) (0.4242,-0.0097743) (0.44962,-0.0075688) (0.4752,-0.0053929) (0.50087,-0.0032675) (0.52657,-0.0012115) (0.55221,0.00075811) (0.57756,0.0025658) (0.60275,0.0041574) (0.62773,0.0055267) (0.65241,0.0066711) (0.67674,0.0075916) (0.70065,0.0082928) (0.72407,0.0087821) (0.74693,0.0090702) (0.76918,0.0091701) (0.79074,0.0090975) (0.81154,0.0088696) (0.83154,0.0085052) (0.85067,0.0080245) (0.86887,0.0074483) (0.88608,0.0067976) (0.90225,0.0060937) (0.91734,0.0053575) (0.93128,0.0046091) (0.94404,0.0038681) (0.95557,0.0031526) (0.96584,0.0024794) (0.97481,0.0018638) (0.98245,0.0013192) (0.98874,0.00085744) (0.99365,0.00048808) (0.99718,0.00021876) (0.99929,5.4967e-05) (1,1.6257e-17)};\addlegendentry{$p=0{,}55$}
\addplot[corange,dashed,thin,forget plot] coordinates {(0,0) (0.016949,0.0030342) (0.033898,0.0059734) (0.050847,0.0088176) (0.067797,0.011567) (0.084746,0.014221) (0.10169,0.016781) (0.11864,0.019245) (0.13559,0.021614) (0.15254,0.023889) (0.16949,0.026068) (0.18644,0.028153) (0.20339,0.030142) (0.22034,0.032037) (0.23729,0.033837) (0.25424,0.035541) (0.27119,0.037151) (0.28814,0.038666) (0.30508,0.040085) (0.32203,0.04141) (0.33898,0.04264) (0.35593,0.043775) (0.37288,0.044815) (0.38983,0.04576) (0.40678,0.04661) (0.42373,0.047365) (0.44068,0.048025) (0.45763,0.04859) (0.47458,0.04906) (0.49153,0.049435) (0.50847,0.049715) (0.52542,0.0499) (0.54237,0.04999) (0.55932,0.049979) (0.57627,0.04983) (0.59322,0.049539) (0.61017,0.049106) (0.62712,0.048532) (0.64407,0.047815) (0.66102,0.046957) (0.67797,0.045957) (0.69492,0.044815) (0.71186,0.043531) (0.72881,0.042105) (0.74576,0.040538) (0.76271,0.038828) (0.77966,0.036977) (0.79661,0.034984) (0.81356,0.032849) (0.83051,0.030572) (0.84746,0.028153) (0.86441,0.025592) (0.88136,0.02289) (0.89831,0.020045) (0.91525,0.017059) (0.9322,0.013931) (0.94915,0.010661) (0.9661,0.0072492) (0.98305,0.0036955) (1,0)};
\nextgroupplot[title={effet de $t$ ($m=2$~\%, $p=0{,}4$)}]
\addplot[cmuted,thick] coordinates {(1,-8.3082e-18) (0.99932,9.5323e-05) (0.99727,0.00038046) (0.99387,0.00085291) (0.98913,0.0015086) (0.98304,0.0023417) (0.97564,0.0033453) (0.96694,0.0045107) (0.95696,0.0058281) (0.94573,0.0072865) (0.93329,0.0088741) (0.91966,0.010578) (0.90487,0.012385) (0.88898,0.01428) (0.87202,0.01625) (0.85404,0.018279) (0.83509,0.020352) (0.81521,0.022454) (0.79447,0.02457) (0.77291,0.026685) (0.75061,0.028785) (0.72761,0.030854) (0.70397,0.032878) (0.67978,0.034843) (0.65508,0.036736) (0.62996,0.038544) (0.60447,0.040253) (0.57868,0.041852) (0.55268,0.043329) (0.52653,0.044673) (0.50029,0.045874) (0.47406,0.046922) (0.44788,0.047809) (0.42185,0.048528) (0.39602,0.04907) (0.37037,0.049374) (0.34509,0.049404) (0.32022,0.049164) (0.29586,0.048657) (0.27206,0.04789) (0.24889,0.04687) (0.22642,0.04561) (0.20471,0.044123) (0.18383,0.042422) (0.16383,0.040526) (0.14477,0.038452) (0.1267,0.03622) (0.10968,0.033851) (0.093744,0.031367) (0.078945,0.02879) (0.065322,0.026141) (0.05291,0.023443) (0.041744,0.020718) (0.031853,0.017984) (0.023262,0.015262) (0.015993,0.012569) (0.010066,0.0099208) (0.0054943,0.00733) (0.0022887,0.0048079) (0.0004562,0.0023627) (0,0) (0.00091426,-0.0022258) (0.0031894,-0.004262) (0.0068174,-0.0061083) (0.011786,-0.0077654) (0.018081,-0.0092345) (0.025682,-0.010518) (0.034567,-0.011618) (0.04471,-0.012539) (0.056083,-0.013286) (0.068653,-0.013866) (0.082384,-0.014284) (0.097239,-0.014549) (0.11318,-0.01467) (0.13015,-0.014658) (0.14813,-0.014524) (0.16704,-0.014281) (0.18685,-0.013942) (0.2075,-0.013521) (0.22894,-0.013034) (0.25111,-0.012495) (0.27395,-0.011921) (0.29741,-0.011328) (0.32141,-0.010732) (0.3459,-0.010147) (0.37081,-0.0095897) (0.39607,-0.009074) (0.42171,-0.0085804) (0.44759,-0.0080623) (0.47361,-0.0075275) (0.49971,-0.0069847) (0.52581,-0.0064414) (0.55185,-0.0059048) (0.57775,-0.005381) (0.60344,-0.0048751) (0.62886,-0.0043915) (0.65393,-0.0039336) (0.67859,-0.0035038) (0.70276,-0.0031038) (0.72639,-0.0027345) (0.74939,-0.002396) (0.77172,-0.0020879) (0.79332,-0.0018092) (0.81411,-0.0015587) (0.83404,-0.0013347) (0.85306,-0.0011353) (0.87112,-0.00095864) (0.88816,-0.00080276) (0.90414,-0.00066572) (0.91901,-0.00054571) (0.93274,-0.00044107) (0.94527,-0.00035028) (0.95658,-0.00027204) (0.96664,-0.00020524) (0.97542,-0.00014893) (0.98288,-0.00010238) (0.98902,-6.5003e-05) (0.99381,-3.6343e-05) (0.99725,-1.6084e-05) (0.99931,-4.011e-06) (1,8.3082e-18)};\addlegendentry{$t=6~\%$}
\addplot[cblue,very thick] coordinates {(1,-1.6616e-17) (0.99932,0.00014499) (0.99729,0.00057873) (0.9939,0.0012975) (0.98918,0.0022953) (0.98312,0.0035638) (0.97575,0.0050924) (0.96709,0.0068686) (0.95715,0.0088781) (0.94597,0.011105) (0.93356,0.013532) (0.91998,0.01614) (0.90524,0.01891) (0.88939,0.021822) (0.87247,0.024854) (0.85453,0.027985) (0.83561,0.031195) (0.81577,0.03446) (0.79505,0.03776) (0.77351,0.041072) (0.75121,0.044375) (0.72822,0.047648) (0.70458,0.050869) (0.68037,0.054017) (0.65566,0.057071) (0.6305,0.060012) (0.60498,0.062817) (0.57915,0.065468) (0.5531,0.067946) (0.52688,0.07023) (0.50059,0.072303) (0.47428,0.074147) (0.44803,0.075745) (0.42192,0.077082) (0.39599,0.078142) (0.37016,0.078855) (0.34468,0.07918) (0.31963,0.079112) (0.29508,0.07865) (0.27111,0.077795) (0.24777,0.076553) (0.22515,0.074933) (0.20331,0.072945) (0.18232,0.070605) (0.16222,0.06793) (0.14309,0.06494) (0.12497,0.061659) (0.10793,0.058112) (0.091996,0.054325) (0.077226,0.050327) (0.063657,0.046145) (0.051324,0.041808) (0.040261,0.037346) (0.030495,0.032785) (0.022051,0.028152) (0.01495,0.023471) (0.0092059,0.018764) (0.0048327,0.014049) (0.0018383,0.0093428) (0.00022717,0.0046569) (0,0) (0.0011433,-0.00452) (0.0036398,-0.0087969) (0.007479,-0.012828) (0.012647,-0.016609) (0.019125,-0.020136) (0.026892,-0.023408) (0.035924,-0.026419) (0.046194,-0.029168) (0.057669,-0.031651) (0.070318,-0.033869) (0.084103,-0.03582) (0.098987,-0.037507) (0.11493,-0.038931) (0.13188,-0.040097) (0.1498,-0.041012) (0.16865,-0.041685) (0.18836,-0.042124) (0.2089,-0.042343) (0.23021,-0.042356) (0.25223,-0.042178) (0.2749,-0.041827) (0.29818,-0.041321) (0.322,-0.040679) (0.3463,-0.039923) (0.37102,-0.039071) (0.3961,-0.038146) (0.42164,-0.037134) (0.44744,-0.035998) (0.47339,-0.034752) (0.49941,-0.033414) (0.52545,-0.031999) (0.55143,-0.030522) (0.57728,-0.028997) (0.60293,-0.027439) (0.62832,-0.025859) (0.65336,-0.024269) (0.67799,-0.022677) (0.70216,-0.021095) (0.72577,-0.019529) (0.74879,-0.017986) (0.77113,-0.016475) (0.79274,-0.014999) (0.81355,-0.013565) (0.83352,-0.012178) (0.85258,-0.010842) (0.87067,-0.0095628) (0.88775,-0.0083442) (0.90378,-0.0071909) (0.91869,-0.0061076) (0.93246,-0.0050987) (0.94504,-0.0041687) (0.9564,-0.0033221) (0.96649,-0.0025632) (0.9753,-0.001896) (0.9828,-0.0013244) (0.98897,-0.00085179) (0.99379,-0.00048097) (0.99723,-0.00021436) (0.99931,-5.3678e-05) (1,1.6616e-17)};\addlegendentry{$t=12~\%$}
\addplot[corange,very thick] coordinates {(1,-2.7694e-17) (0.99933,0.00021121) (0.9973,0.00084309) (0.99394,0.0018904) (0.98925,0.0033444) (0.98323,0.0051932) (0.9759,0.0074219) (0.96729,0.010013) (0.9574,0.012945) (0.94627,0.016196) (0.93393,0.019742) (0.92041,0.023556) (0.90573,0.02761) (0.88994,0.031877) (0.87308,0.036326) (0.85518,0.040928) (0.83631,0.045652) (0.8165,0.050469) (0.79582,0.055346) (0.7743,0.060254) (0.75202,0.065162) (0.72903,0.07004) (0.70539,0.074857) (0.68117,0.079582) (0.65643,0.084185) (0.63123,0.088635) (0.60566,0.092903) (0.57978,0.096957) (0.55365,0.10077) (0.52736,0.10431) (0.50098,0.10754) (0.47458,0.11045) (0.44823,0.11299) (0.42201,0.11515) (0.39595,0.1169) (0.36987,0.11816) (0.34414,0.11888) (0.31884,0.11904) (0.29405,0.11864) (0.26984,0.11767) (0.24629,0.11613) (0.22347,0.11403) (0.20145,0.11137) (0.1803,0.10818) (0.16008,0.10447) (0.14085,0.10026) (0.12267,0.095578) (0.10559,0.09046) (0.089666,0.084936) (0.074934,0.079042) (0.061436,0.072816) (0.049209,0.066295) (0.038283,0.059517) (0.028686,0.05252) (0.020438,0.045339) (0.013558,0.038007) (0.008059,0.030555) (0.0039506,0.023008) (0.0012378,0.015389) (-7.8199e-05,0.0077159) (0,0) (0.0014487,-0.007579) (0.0042403,-0.014843) (0.0083611,-0.021786) (0.013793,-0.028399) (0.020516,-0.034672) (0.028506,-0.040594) (0.037734,-0.046154) (0.048171,-0.051339) (0.059784,-0.056138) (0.072538,-0.06054) (0.086396,-0.064536) (0.10132,-0.068117) (0.11726,-0.071278) (0.13418,-0.074016) (0.15204,-0.07633) (0.17079,-0.078223) (0.19038,-0.079701) (0.21076,-0.080773) (0.23189,-0.081452) (0.25371,-0.081755) (0.27617,-0.081701) (0.29922,-0.081311) (0.32279,-0.08061) (0.34684,-0.079623) (0.37131,-0.07838) (0.39614,-0.076909) (0.42155,-0.075206) (0.44724,-0.073245) (0.47309,-0.071052) (0.49902,-0.068653) (0.52497,-0.066075) (0.55088,-0.063344) (0.57666,-0.060486) (0.60225,-0.057525) (0.62759,-0.054483) (0.65259,-0.051382) (0.6772,-0.048242) (0.70135,-0.045082) (0.72496,-0.041921) (0.74798,-0.038774) (0.77034,-0.035657) (0.79197,-0.032585) (0.81282,-0.029573) (0.83282,-0.026635) (0.85192,-0.023785) (0.87007,-0.021035) (0.88721,-0.018399) (0.90329,-0.015891) (0.91827,-0.013523) (0.93209,-0.011309) (0.94473,-0.0092599) (0.95614,-0.0073888) (0.9663,-0.0057071) (0.97516,-0.0042255) (0.9827,-0.0029539) (0.9889,-0.0019008) (0.99375,-0.0010738) (0.99722,-0.00047872) (0.9993,-0.0001199) (1,2.7694e-17)};\addlegendentry{$t=20~\%$}
\end{groupplot}
\end{tikzpicture}

\caption{Effet de chaque paramètre (calculé par \code{naca4}). Tirets : ligne moyenne.}\label{fig:parametres}
\end{figure}
\begin{itemize}
\item \textbf{L'épaisseur $t$} pilote le rayon de bord d'attaque ($r_{BA}\propto t^2$), donc la tolérance en incidence et la forme du pic d'aspiration. Elle fixe aussi la survitesse moyenne sur le profil, donc le gradient de pression adverse en aval et la traînée de forme. Elle augmente légèrement la pente de portance, $\dd C_L/\dd\alpha\approx2\pi(1+0{,}77\,t)$.
\item \textbf{La cambrure $m$} décale la courbe de portance. En théorie des profils minces, $C_L=2\pi(\alpha-\alpha_{L0})$ avec un angle de portance nulle $\alpha_{L0}<0$ proportionnel à $m$. Elle crée un moment piqueur, $C_{M,c/4}<0$.
\item \textbf{La position $p$} répartit la charge le long de la corde. Une cambrure arrière ($p$ grand) charge le bord de fuite, augmente $|C_{M}|$ et crée un gradient adverse plus fort en fin d'extrados. Une cambrure avant renforce le pic d'aspiration au bord d'attaque.
\end{itemize}

\begin{figure}[htbp]
\centering
% figure générée par tools/make_figures.py : profils
\begin{tikzpicture}
\begin{axis}[width=\linewidth,height=4.6cm,axis equal image,xmin=-0.02,xmax=1.02,ymin=-0.12,ymax=0.16,
  xlabel={$x/c$},ylabel={$y/c$},legend columns=3,legend style={draw=none,at={(0.5,1.03)},anchor=south,font=\small}]
\addplot[cmuted,thick] coordinates {(1,-1.6653e-17) (0.99961,5.6016e-05) (0.99846,0.00022381) (0.99653,0.00050262) (0.99384,0.00089119) (0.99039,0.0013878) (0.98618,0.0019902) (0.98123,0.0026957) (0.97553,0.0035014) (0.9691,0.0044037) (0.96194,0.0053988) (0.95407,0.0064826) (0.9455,0.0076508) (0.93625,0.0088987) (0.92632,0.010221) (0.91573,0.011614) (0.90451,0.013071) (0.89266,0.014587) (0.8802,0.016158) (0.86716,0.017777) (0.85355,0.019438) (0.8394,0.021138) (0.82472,0.022869) (0.80955,0.024626) (0.79389,0.026405) (0.77779,0.028198) (0.76125,0.03) (0.74431,0.031806) (0.727,0.03361) (0.70933,0.035406) (0.69134,0.037188) (0.67306,0.03895) (0.65451,0.040686) (0.63572,0.04239) (0.61672,0.044055) (0.59755,0.045674) (0.57822,0.047242) (0.55877,0.048752) (0.53923,0.050196) (0.51963,0.051568) (0.5,0.052862) (0.48037,0.054069) (0.46077,0.055184) (0.44123,0.056199) (0.42178,0.057108) (0.40245,0.057905) (0.38328,0.058582) (0.36428,0.059135) (0.34549,0.059557) (0.32694,0.059843) (0.30866,0.059988) (0.29067,0.059989) (0.273,0.059841) (0.25569,0.059542) (0.23875,0.059088) (0.22221,0.058478) (0.20611,0.057712) (0.19045,0.056788) (0.17528,0.055708) (0.1606,0.054472) (0.14645,0.053083) (0.13284,0.051542) (0.1198,0.049854) (0.10734,0.048021) (0.095492,0.046049) (0.084265,0.043942) (0.07368,0.041705) (0.063752,0.039345) (0.054497,0.036867) (0.045928,0.034276) (0.03806,0.03158) (0.030904,0.028783) (0.024472,0.025893) (0.018772,0.022915) (0.013815,0.019854) (0.0096074,0.016715) (0.0061558,0.013503) (0.0034658,0.010223) (0.0015413,0.0068767) (0.00038548,0.0034684) (0,0) (0.00038548,-0.0034684) (0.0015413,-0.0068767) (0.0034658,-0.010223) (0.0061558,-0.013503) (0.0096074,-0.016715) (0.013815,-0.019854) (0.018772,-0.022915) (0.024472,-0.025893) (0.030904,-0.028783) (0.03806,-0.03158) (0.045928,-0.034276) (0.054497,-0.036867) (0.063752,-0.039345) (0.07368,-0.041705) (0.084265,-0.043942) (0.095492,-0.046049) (0.10734,-0.048021) (0.1198,-0.049854) (0.13284,-0.051542) (0.14645,-0.053083) (0.1606,-0.054472) (0.17528,-0.055708) (0.19045,-0.056788) (0.20611,-0.057712) (0.22221,-0.058478) (0.23875,-0.059088) (0.25569,-0.059542) (0.273,-0.059841) (0.29067,-0.059989) (0.30866,-0.059988) (0.32694,-0.059843) (0.34549,-0.059557) (0.36428,-0.059135) (0.38328,-0.058582) (0.40245,-0.057905) (0.42178,-0.057108) (0.44123,-0.056199) (0.46077,-0.055184) (0.48037,-0.054069) (0.5,-0.052862) (0.51963,-0.051568) (0.53923,-0.050196) (0.55877,-0.048752) (0.57822,-0.047242) (0.59755,-0.045674) (0.61672,-0.044055) (0.63572,-0.04239) (0.65451,-0.040686) (0.67306,-0.03895) (0.69134,-0.037188) (0.70933,-0.035406) (0.727,-0.03361) (0.74431,-0.031806) (0.76125,-0.03) (0.77779,-0.028198) (0.79389,-0.026405) (0.80955,-0.024626) (0.82472,-0.022869) (0.8394,-0.021138) (0.85355,-0.019438) (0.86716,-0.017777) (0.8802,-0.016158) (0.89266,-0.014587) (0.90451,-0.013071) (0.91573,-0.011614) (0.92632,-0.010221) (0.93625,-0.0088987) (0.9455,-0.0076508) (0.95407,-0.0064826) (0.96194,-0.0053988) (0.9691,-0.0044037) (0.97553,-0.0035014) (0.98123,-0.0026957) (0.98618,-0.0019902) (0.99039,-0.0013878) (0.99384,-0.00089119) (0.99653,-0.00050262) (0.99846,-0.00022381) (0.99961,-5.6016e-05) (1,1.6653e-17)};
\addlegendentry{NACA 0012}
\addplot[cblue,very thick] coordinates {(1,-1.6616e-17) (0.99962,8.1582e-05) (0.99847,0.00032594) (0.99657,0.0007319) (0.9939,0.0012975) (0.99048,0.0020202) (0.98631,0.0028964) (0.9814,0.003922) (0.97575,0.0050924) (0.96937,0.0064021) (0.96228,0.0078452) (0.95447,0.0094151) (0.94597,0.011105) (0.93678,0.012907) (0.92692,0.014814) (0.9164,0.016818) (0.90524,0.01891) (0.89346,0.021082) (0.88106,0.023324) (0.86808,0.025628) (0.85453,0.027985) (0.84043,0.030386) (0.8258,0.032822) (0.81067,0.035283) (0.79505,0.03776) (0.77897,0.040244) (0.76245,0.042726) (0.74553,0.045197) (0.72822,0.047648) (0.71055,0.05007) (0.69254,0.052453) (0.67424,0.05479) (0.65566,0.057071) (0.63683,0.059288) (0.61778,0.061432) (0.59855,0.063495) (0.57915,0.065468) (0.55963,0.067344) (0.54001,0.069113) (0.52032,0.070769) (0.50059,0.072303) (0.48085,0.073708) (0.46114,0.074977) (0.44149,0.076104) (0.42192,0.077082) (0.40247,0.077904) (0.38303,0.078547) (0.36375,0.078973) (0.34468,0.07918) (0.32585,0.079166) (0.30729,0.07893) (0.28903,0.078473) (0.27111,0.077795) (0.25354,0.0769) (0.23637,0.07579) (0.21962,0.07447) (0.20331,0.072945) (0.18748,0.071222) (0.17215,0.069308) (0.15735,0.067211) (0.14309,0.06494) (0.1294,0.062506) (0.11631,0.059917) (0.10384,0.057187) (0.091996,0.054325) (0.080807,0.051345) (0.070289,0.048257) (0.060456,0.045074) (0.051324,0.041808) (0.042906,0.038472) (0.035214,0.035076) (0.02826,0.031633) (0.022051,0.028152) (0.016598,0.024645) (0.011907,0.02112) (0.0079837,0.017585) (0.0048327,0.014049) (0.0024573,0.010518) (0.00085968,0.0069967) (4.0695e-05,0.0034897) (0,0) (0.00073027,-0.0034127) (0.002223,-0.006689) (0.0044742,-0.0098278) (0.007479,-0.012828) (0.011231,-0.015687) (0.015723,-0.018404) (0.020946,-0.020978) (0.026892,-0.023408) (0.033549,-0.025691) (0.040906,-0.027826) (0.048951,-0.029813) (0.057669,-0.031651) (0.067048,-0.03334) (0.077071,-0.034878) (0.087723,-0.036267) (0.098987,-0.037507) (0.11085,-0.038599) (0.12328,-0.039546) (0.13627,-0.040349) (0.1498,-0.041012) (0.16385,-0.041539) (0.1784,-0.041933) (0.19342,-0.042199) (0.2089,-0.042343) (0.22481,-0.042372) (0.24113,-0.04229) (0.25784,-0.042106) (0.2749,-0.041827) (0.29231,-0.041461) (0.31003,-0.041016) (0.32803,-0.0405) (0.3463,-0.039923) (0.36481,-0.039292) (0.38352,-0.038617) (0.40244,-0.037905) (0.42164,-0.037134) (0.44097,-0.036293) (0.4604,-0.035388) (0.47989,-0.034426) (0.49941,-0.033414) (0.51894,-0.032359) (0.53845,-0.031267) (0.55791,-0.030144) (0.57728,-0.028997) (0.59654,-0.027831) (0.61566,-0.026651) (0.63461,-0.025462) (0.65336,-0.024269) (0.67188,-0.023075) (0.69014,-0.021884) (0.70811,-0.020701) (0.72577,-0.019529) (0.74309,-0.018369) (0.76005,-0.017226) (0.7766,-0.016102) (0.79274,-0.014999) (0.80843,-0.013919) (0.82365,-0.012865) (0.83837,-0.011839) (0.85258,-0.010842) (0.86624,-0.0098771) (0.87934,-0.0089456) (0.89186,-0.0080496) (0.90378,-0.0071909) (0.91507,-0.0063716) (0.92572,-0.0055935) (0.93572,-0.0048586) (0.94504,-0.0041687) (0.95367,-0.0035257) (0.9616,-0.0029314) (0.96882,-0.0023876) (0.9753,-0.001896) (0.98105,-0.0014582) (0.98606,-0.0010756) (0.9903,-0.00074943) (0.99379,-0.00048097) (0.9965,-0.00027113) (0.99844,-0.00012069) (0.99961,-3.0201e-05) (1,1.6616e-17)};
\addlegendentry{NACA 2412}
\addplot[corange,very thick] coordinates {(1,-2.0634e-17) (0.99962,0.00012079) (0.9985,0.00048257) (0.99662,0.0010836) (0.99399,0.001921) (0.99062,0.0029907) (0.98651,0.0042877) (0.98166,0.0058057) (0.97608,0.0075377) (0.96979,0.0094755) (0.96278,0.01161) (0.95506,0.013932) (0.94665,0.01643) (0.93756,0.019094) (0.9278,0.021911) (0.91739,0.024869) (0.90633,0.027956) (0.89464,0.031158) (0.88235,0.034462) (0.86946,0.037853) (0.85599,0.041319) (0.84197,0.044845) (0.82741,0.048416) (0.81234,0.05202) (0.79677,0.055641) (0.78073,0.059266) (0.76425,0.06288) (0.74734,0.06647) (0.73004,0.070022) (0.71236,0.073522) (0.69435,0.076957) (0.67601,0.080314) (0.65738,0.083579) (0.63849,0.086741) (0.61937,0.089786) (0.60005,0.092702) (0.58055,0.095477) (0.56092,0.098101) (0.54117,0.10056) (0.52134,0.10285) (0.50147,0.10495) (0.48158,0.10686) (0.4617,0.10856) (0.44187,0.11006) (0.42213,0.11133) (0.40249,0.11238) (0.38267,0.11316) (0.36296,0.11359) (0.34346,0.11368) (0.32421,0.11342) (0.30524,0.11282) (0.28658,0.11189) (0.26826,0.11062) (0.25033,0.10903) (0.23282,0.10712) (0.21574,0.10491) (0.19915,0.1024) (0.18306,0.099621) (0.1675,0.096574) (0.15251,0.093279) (0.1381,0.089754) (0.12431,0.086017) (0.11115,0.082086) (0.09865,0.077982) (0.086827,0.073724) (0.0757,0.069333) (0.065285,0.06483) (0.055598,0.060235) (0.046652,0.055567) (0.038459,0.050847) (0.031031,0.046094) (0.024375,0.041324) (0.018499,0.036555) (0.013409,0.031803) (0.0091099,0.027083) (0.0056045,0.022405) (0.0028946,0.017783) (0.00098061,0.013225) (-0.00013821,0.0087379) (-0.00046398,0.0043285) (0,0) (0.0012349,-0.0041744) (0.0032209,-0.0081225) (0.0059509,-0.011844) (0.0094171,-0.015339) (0.01361,-0.018609) (0.01852,-0.021652) (0.024136,-0.024471) (0.030445,-0.027066) (0.037434,-0.02944) (0.04509,-0.031594) (0.053397,-0.033531) (0.062341,-0.035254) (0.071906,-0.036766) (0.082075,-0.038073) (0.09283,-0.039178) (0.10416,-0.040087) (0.11603,-0.040806) (0.12844,-0.041343) (0.14137,-0.041704) (0.15479,-0.041899) (0.16869,-0.041936) (0.18305,-0.041825) (0.19785,-0.041576) (0.21307,-0.041202) (0.22869,-0.040713) (0.24469,-0.040121) (0.26105,-0.03944) (0.27774,-0.038683) (0.29476,-0.037863) (0.31208,-0.036993) (0.32967,-0.036088) (0.34752,-0.035161) (0.3656,-0.034226) (0.38389,-0.033295) (0.40242,-0.032382) (0.42144,-0.031437) (0.44059,-0.030435) (0.45984,-0.029384) (0.47916,-0.028293) (0.49853,-0.027172) (0.51792,-0.026028) (0.53729,-0.024869) (0.55662,-0.023702) (0.57588,-0.022535) (0.59504,-0.021374) (0.61407,-0.020223) (0.63295,-0.019088) (0.65164,-0.017974) (0.67011,-0.016883) (0.68834,-0.015819) (0.70629,-0.014785) (0.72395,-0.013783) (0.74128,-0.012814) (0.75825,-0.01188) (0.77484,-0.010982) (0.79101,-0.010119) (0.80676,-0.0092929) (0.82204,-0.0085032) (0.83683,-0.0077497) (0.85112,-0.0070324) (0.86487,-0.0063508) (0.87806,-0.0057048) (0.89067,-0.005094) (0.90269,-0.0045179) (0.91408,-0.0039765) (0.92484,-0.0034694) (0.93493,-0.0029965) (0.94435,-0.0025578) (0.95308,-0.0021532) (0.9611,-0.0017827) (0.96841,-0.0014466) (0.97497,-0.001145) (0.9808,-0.00087806) (0.98586,-0.00064606) (0.99017,-0.00044925) (0.9937,-0.00028784) (0.99645,-0.00016206) (0.99842,-7.2072e-05) (0.99961,-1.8025e-05) (1,2.0634e-17)};
\addlegendentry{NACA 4415}
\end{axis}
\end{tikzpicture}

\caption{Trois profils de référence du cours, tracés à partir des points du code.}\label{fig:profils}
\end{figure}

\begin{piege}
Pour $m=0$, la position $p$ n'a aucun effet. Sur la grille $8\times8\times8$, les 8 valeurs de $p$ à $m=0$ donnent le même profil symétrique. Ce sont des doublons dans le front de Pareto, et des minima dégénérés dans le QUBO. Le code gère ce cas en renvoyant $y_c=0$ dès que $m=0$ ou $p=0$, ce qui évite aussi une division par zéro.
\end{piege}

\section{La discrétisation : distribution en cosinus}
Le contour est échantillonné en $2n+1$ points ($n=$\,\code{n\_half}, 100 par défaut, soit 200 panneaux). Les abscisses suivent une \textbf{distribution en cosinus} :
\[
x_k=\tfrac12\big(1-\cos\beta_k\big),\qquad \beta_k=\frac{k\pi}{n},\quad k=0,\dots,n .
\]
C'est la projection sur l'axe des $x$ de points également espacés sur un demi-cercle. Les points se resserrent aux deux bords, où $\dd x/\dd\beta=\tfrac12\sin\beta\to0$. Le premier intervalle vaut $\tfrac12(1-\cos\tfrac\pi n)\approx\pi^2/(4n^2)$, soit $2{,}5\cdot10^{-4}$ pour $n=100$, contre $10^{-2}$ en distribution uniforme : quarante fois plus fin.

C'est indispensable :
\begin{itemize}
\item au \textbf{bord d'attaque}, la courbure est maximale, la vitesse passe de 0 (point d'arrêt) à son maximum sur quelques millièmes de corde, et la couche limite démarre là ;
\item au \textbf{bord de fuite}, la condition de Kutta porte sur les deux derniers panneaux, et la couche limite y est la plus sensible.
\end{itemize}

\begin{figure}[htbp]
\centering
% figure générée par tools/make_figures.py : cosinus
\begin{tikzpicture}
\begin{axis}[width=0.62\linewidth,height=5cm,xmin=-0.005,xmax=0.2,ymin=0,ymax=0.065,scaled ticks=false,
  xtick={0,0.05,0.1,0.15,0.2},ytick={0,0.02,0.04,0.06},tick label style={/pgf/number format/fixed,/pgf/number format/use comma,font=\small},
  xlabel={$x/c$},ylabel={$y/c$},legend style={draw=none,font=\small,at={(0.98,0.04)},anchor=south east}]
\addplot[cmuted,thin,forget plot] coordinates {(0,0) (0.0025063,0.0087274) (0.0050125,0.012228) (0.0075188,0.014866) (0.010025,0.017057) (0.012531,0.018961) (0.015038,0.020661) (0.017544,0.022205) (0.02005,0.023625) (0.022556,0.024944) (0.025063,0.026177) (0.027569,0.027337) (0.030075,0.028433) (0.032581,0.029474) (0.035088,0.030464) (0.037594,0.031409) (0.0401,0.032313) (0.042607,0.033179) (0.045113,0.034012) (0.047619,0.034813) (0.050125,0.035585) (0.052632,0.036329) (0.055138,0.037048) (0.057644,0.037743) (0.06015,0.038416) (0.062657,0.039067) (0.065163,0.039698) (0.067669,0.04031) (0.070175,0.040904) (0.072682,0.04148) (0.075188,0.04204) (0.077694,0.042585) (0.080201,0.043114) (0.082707,0.043629) (0.085213,0.04413) (0.087719,0.044617) (0.090226,0.045092) (0.092732,0.045554) (0.095238,0.046004) (0.097744,0.046443) (0.10025,0.04687) (0.10276,0.047286) (0.10526,0.047692) (0.10777,0.048088) (0.11028,0.048474) (0.11278,0.04885) (0.11529,0.049217) (0.11779,0.049574) (0.1203,0.049923) (0.12281,0.050263) (0.12531,0.050595) (0.12782,0.050918) (0.13033,0.051234) (0.13283,0.051541) (0.13534,0.051841) (0.13784,0.052134) (0.14035,0.052419) (0.14286,0.052697) (0.14536,0.052968) (0.14787,0.053232) (0.15038,0.053489) (0.15288,0.05374) (0.15539,0.053984) (0.15789,0.054222) (0.1604,0.054454) (0.16291,0.05468) (0.16541,0.054899) (0.16792,0.055113) (0.17043,0.055321) (0.17293,0.055524) (0.17544,0.055721) (0.17794,0.055912) (0.18045,0.056098) (0.18296,0.056279) (0.18546,0.056454) (0.18797,0.056624) (0.19048,0.05679) (0.19298,0.05695) (0.19549,0.057106) (0.19799,0.057256) (0.2005,0.057402) (0.20301,0.057543) (0.20551,0.05768) (0.20802,0.057812) (0.21053,0.05794) (0.21303,0.058063) (0.21554,0.058182) (0.21805,0.058297) (0.22055,0.058407) (0.22306,0.058514) (0.22556,0.058616) (0.22807,0.058714) (0.23058,0.058808) (0.23308,0.058898) (0.23559,0.058984) (0.2381,0.059067) (0.2406,0.059146) (0.24311,0.05922) (0.24561,0.059292) (0.24812,0.059359) (0.25063,0.059423) (0.25313,0.059484) (0.25564,0.05954) (0.25815,0.059594) (0.26065,0.059644) (0.26316,0.05969) (0.26566,0.059734) (0.26817,0.059773) (0.27068,0.05981) (0.27318,0.059843) (0.27569,0.059874) (0.2782,0.059901) (0.2807,0.059924) (0.28321,0.059945) (0.28571,0.059963) (0.28822,0.059978) (0.29073,0.059989) (0.29323,0.059998) (0.29574,0.060004) (0.29825,0.060007) (0.30075,0.060007) (0.30326,0.060004) (0.30576,0.059998) (0.30827,0.05999) (0.31078,0.059979) (0.31328,0.059965) (0.31579,0.059948) (0.3183,0.059929) (0.3208,0.059907) (0.32331,0.059883) (0.32581,0.059856) (0.32832,0.059826) (0.33083,0.059794) (0.33333,0.05976) (0.33584,0.059722) (0.33835,0.059683) (0.34085,0.059641) (0.34336,0.059596) (0.34586,0.05955) (0.34837,0.0595) (0.35088,0.059449) (0.35338,0.059395) (0.35589,0.059339) (0.3584,0.059281) (0.3609,0.05922) (0.36341,0.059157) (0.36591,0.059092) (0.36842,0.059025) (0.37093,0.058955) (0.37343,0.058884) (0.37594,0.05881) (0.37845,0.058734) (0.38095,0.058656) (0.38346,0.058576) (0.38596,0.058494) (0.38847,0.05841) (0.39098,0.058324) (0.39348,0.058236) (0.39599,0.058146) (0.3985,0.058054) (0.401,0.05796) (0.40351,0.057864) (0.40602,0.057766) (0.40852,0.057667) (0.41103,0.057565) (0.41353,0.057462) (0.41604,0.057356) (0.41855,0.057249) (0.42105,0.05714) (0.42356,0.05703) (0.42607,0.056917) (0.42857,0.056803) (0.43108,0.056687) (0.43358,0.056569) (0.43609,0.05645) (0.4386,0.056328) (0.4411,0.056205) (0.44361,0.056081) (0.44612,0.055955) (0.44862,0.055827) (0.45113,0.055697) (0.45363,0.055566) (0.45614,0.055433) (0.45865,0.055299) (0.46115,0.055163) (0.46366,0.055025) (0.46617,0.054886) (0.46867,0.054746) (0.47118,0.054603) (0.47368,0.05446) (0.47619,0.054315) (0.4787,0.054168) (0.4812,0.05402) (0.48371,0.05387) (0.48622,0.053719) (0.48872,0.053566) (0.49123,0.053412) (0.49373,0.053256) (0.49624,0.0531) (0.49875,0.052941) (0.50125,0.052781) (0.50376,0.05262) (0.50627,0.052458) (0.50877,0.052294) (0.51128,0.052129) (0.51378,0.051962) (0.51629,0.051794) (0.5188,0.051625) (0.5213,0.051454) (0.52381,0.051282) (0.52632,0.051109) (0.52882,0.050934) (0.53133,0.050759) (0.53383,0.050581) (0.53634,0.050403) (0.53885,0.050223) (0.54135,0.050043) (0.54386,0.04986) (0.54637,0.049677) (0.54887,0.049492) (0.55138,0.049307) (0.55388,0.04912) (0.55639,0.048931) (0.5589,0.048742) (0.5614,0.048551) (0.56391,0.048359) (0.56642,0.048166) (0.56892,0.047972) (0.57143,0.047777) (0.57393,0.047581) (0.57644,0.047383) (0.57895,0.047184) (0.58145,0.046984) (0.58396,0.046783) (0.58647,0.046581) (0.58897,0.046378) (0.59148,0.046173) (0.59398,0.045968) (0.59649,0.045761) (0.599,0.045554) (0.6015,0.045345) (0.60401,0.045135) (0.60652,0.044924) (0.60902,0.044712) (0.61153,0.044499) (0.61404,0.044285) (0.61654,0.04407) (0.61905,0.043854) (0.62155,0.043637) (0.62406,0.043419) (0.62657,0.043199) (0.62907,0.042979) (0.63158,0.042758) (0.63409,0.042535) (0.63659,0.042312) (0.6391,0.042088) (0.6416,0.041862) (0.64411,0.041636) (0.64662,0.041409) (0.64912,0.04118) (0.65163,0.040951) (0.65414,0.040721) (0.65664,0.040489) (0.65915,0.040257) (0.66165,0.040024) (0.66416,0.03979) (0.66667,0.039554) (0.66917,0.039318) (0.67168,0.039081) (0.67419,0.038843) (0.67669,0.038604) (0.6792,0.038364) (0.6817,0.038123) (0.68421,0.037882) (0.68672,0.037639) (0.68922,0.037395) (0.69173,0.03715) (0.69424,0.036905) (0.69674,0.036658) (0.69925,0.036411) (0.70175,0.036163) (0.70426,0.035913) (0.70677,0.035663) (0.70927,0.035412) (0.71178,0.03516) (0.71429,0.034907) (0.71679,0.034653) (0.7193,0.034398) (0.7218,0.034143) (0.72431,0.033886) (0.72682,0.033629) (0.72932,0.03337) (0.73183,0.033111) (0.73434,0.032851) (0.73684,0.03259) (0.73935,0.032328) (0.74185,0.032065) (0.74436,0.031801) (0.74687,0.031536) (0.74937,0.031271) (0.75188,0.031004) (0.75439,0.030737) (0.75689,0.030469) (0.7594,0.0302) (0.7619,0.02993) (0.76441,0.029659) (0.76692,0.029387) (0.76942,0.029114) (0.77193,0.028841) (0.77444,0.028566) (0.77694,0.028291) (0.77945,0.028014) (0.78195,0.027737) (0.78446,0.027459) (0.78697,0.02718) (0.78947,0.0269) (0.79198,0.02662) (0.79449,0.026338) (0.79699,0.026055) (0.7995,0.025772) (0.80201,0.025488) (0.80451,0.025202) (0.80702,0.024916) (0.80952,0.024629) (0.81203,0.024341) (0.81454,0.024052) (0.81704,0.023763) (0.81955,0.023472) (0.82206,0.02318) (0.82456,0.022888) (0.82707,0.022595) (0.82957,0.0223) (0.83208,0.022005) (0.83459,0.021709) (0.83709,0.021412) (0.8396,0.021114) (0.84211,0.020815) (0.84461,0.020516) (0.84712,0.020215) (0.84962,0.019913) (0.85213,0.019611) (0.85464,0.019307) (0.85714,0.019003) (0.85965,0.018697) (0.86216,0.018391) (0.86466,0.018084) (0.86717,0.017776) (0.86967,0.017467) (0.87218,0.017157) (0.87469,0.016846) (0.87719,0.016534) (0.8797,0.016221) (0.88221,0.015907) (0.88471,0.015592) (0.88722,0.015276) (0.88972,0.014959) (0.89223,0.014642) (0.89474,0.014323) (0.89724,0.014003) (0.89975,0.013683) (0.90226,0.013361) (0.90476,0.013038) (0.90727,0.012715) (0.90977,0.01239) (0.91228,0.012064) (0.91479,0.011738) (0.91729,0.01141) (0.9198,0.011081) (0.92231,0.010752) (0.92481,0.010421) (0.92732,0.010089) (0.92982,0.0097563) (0.93233,0.0094225) (0.93484,0.0090877) (0.93734,0.0087517) (0.93985,0.0084148) (0.94236,0.0080768) (0.94486,0.0077377) (0.94737,0.0073976) (0.94987,0.0070564) (0.95238,0.0067141) (0.95489,0.0063708) (0.95739,0.0060264) (0.9599,0.0056808) (0.96241,0.0053342) (0.96491,0.0049865) (0.96742,0.0046377) (0.96992,0.0042878) (0.97243,0.0039367) (0.97494,0.0035846) (0.97744,0.0032313) (0.97995,0.0028768) (0.98246,0.0025213) (0.98496,0.0021646) (0.98747,0.0018067) (0.98997,0.0014477) (0.99248,0.0010875) (0.99499,0.00072621) (0.99749,0.00036369) (1,-1.6653e-17)};
\addplot[only marks,mark=*,mark size=1.6pt,cblue] coordinates {(0,0) (0.0027391,0.0091145) (0.010926,0.01777) (0.024472,0.025893) (0.043227,0.033389) (0.066987,0.040145) (0.095492,0.046049) (0.12843,0.050995) (0.16543,0.054901) (0.20611,0.057712) (0.25,0.059407) (0.29663,0.060005) (0.34549,0.059557) (0.39604,0.058144) (0.44774,0.055872) (0.5,0.052862) (0.55226,0.049241) (0.60396,0.04514) (0.65451,0.040686) (0.70337,0.036002) (0.75,0.031204) (0.79389,0.026405) (0.83457,0.021712) (0.87157,0.017232) (0.90451,0.013071) (0.93301,0.0093315) (0.95677,0.0061118) (0.97553,0.0035014) (0.98907,0.0015769) (0.99726,0.00039741) (1,-1.6653e-17)};
\addlegendentry{cosinus (60 panneaux)}
\addplot[only marks,mark=square*,mark size=1.6pt,corange] coordinates {(0,0) (0.033333,0.029776) (0.066667,0.040067) (0.1,0.046828) (0.13333,0.051602) (0.16667,0.055007) (0.2,0.057373) (0.23333,0.058907) (0.26667,0.05975) (0.3,0.060007) (0.33333,0.05976) (0.36667,0.059072) (0.4,0.057998) (0.43333,0.056581) (0.46667,0.054858) (0.5,0.052862) (0.53333,0.050617) (0.56667,0.048147) (0.6,0.04547) (0.63333,0.042602) (0.66667,0.039554) (0.7,0.036337) (0.73333,0.032955) (0.76667,0.029414) (0.8,0.025715) (0.83333,0.021857) (0.86667,0.017838) (0.9,0.01365) (0.93333,0.0092887) (0.96667,0.0047425) (1,-1.6653e-17)};
\addlegendentry{uniforme (60 panneaux)}
\end{axis}
\end{tikzpicture}

\caption{Zoom sur le bord d'attaque d'un NACA 0012 avec 60 panneaux. La distribution en cosinus place 5 points avant $x=0{,}05$, la distribution uniforme seulement 2 : elle rate complètement l'arrondi du nez.}\label{fig:cosinus}
\end{figure}

\paragraph{Ordre de parcours.} Les points sont rangés comme dans XFOIL : du bord de fuite par l'extrados jusqu'au bord d'attaque, puis par l'intrados jusqu'au bord de fuite. Le contour est donc parcouru dans le sens trigonométrique, et le point de bord de fuite apparaît deux fois (premier et dernier). Avec ce sens, la \textbf{normale sortante} d'un segment de tangente unitaire $(t_x,t_y)$ est $(t_y,-t_x)$ : elle est à droite quand on avance le long du contour. Le même fichier de coordonnées sert d'entrée à XFOIL et au maillage de la SER.

\begin{codebox}[ : \code{aero/naca4.py}]
\begin{lstlisting}
beta = np.linspace(0.0, np.pi, n_half + 1)
x = 0.5 * (1 - np.cos(beta))                 # distribution en cosinus
yt = thickness(x, t, closed_te)              # a4 = -0.1036 : bord de fuite ferme
yc, dyc = camber_line(x, m, p)
th = np.arctan(dyc)
xu, yu = x - yt * np.sin(th), yc + yt * np.cos(th)   # extrados
xl, yl = x + yt * np.sin(th), yc - yt * np.cos(th)   # intrados
X = np.concatenate([xu[::-1], xl[1:]])       # BF extrados -> BA -> BF intrados
Y = np.concatenate([yu[::-1], yl[1:]])
\end{lstlisting}
\code{xu[::-1]} parcourt l'extrados à l'envers (du bord de fuite vers le bord d'attaque), et \code{xl[1:]} évite de dupliquer le point de bord d'attaque. On obtient $X[0]=X[-1]=1$.
\end{codebox}

\begin{projet}
La paramétrisation NACA est pauvre (trois degrés de liberté, des formes toujours « raisonnables ») mais interprétable et énumérable : $8^3=512$ profils. Le rayon de bord d'attaque $r_{BA}=1{,}1019\,t^2$ est exporté dans \code{AeroResult.r\_le}, car c'est probablement le lien physique le plus direct entre les deux objectifs. Il pilote le pic d'aspiration côté aérodynamique, et vraisemblablement le retour radar de face côté SER, ce que l'équipe de la phase~2 doit confirmer. Pour un cas industriel, on passerait à des paramétrisations plus riches (CST de Kulfan, B-splines), donc à plus de variables et plus de qubits.
\end{projet}

\exercices
\begin{exercice}
Pour le NACA 2412 : (a) calculer $y_c(p)$ et la pente de la ligne moyenne en $x=p$ par les deux formules, et vérifier le raccord ; (b) calculer le rayon de bord d'attaque ; (c) que vaut la courbure de la ligne moyenne juste avant et juste après $x=p$ ?
\end{exercice}
\begin{exercice}
Avec $n=100$, comparer l'abscisse du premier point après le bord d'attaque (ou du dernier avant le bord de fuite) en distribution cosinus et en distribution uniforme. Pourquoi le point de bord d'attaque lui-même n'est-il pas un point de collocation ?
\end{exercice}
